\documentclass[12pt]{article}
\usepackage[utf8]{inputenc}
\usepackage[T1]{fontenc}
\usepackage[margin=1in]{geometry}
\usepackage{setspace}
\usepackage{xcolor}
\usepackage[colorlinks=true, urlcolor=blue!60!black, linkcolor=blue!60!black]{hyperref}
\setlength{\parindent}{0em}
\setlength{\parskip}{0.45em}
\renewcommand{\baselinestretch}{1.0}
\begin{document}
\begin{center}
{\Large \textit{Nixon v. Administrator of General Services}, 433 U.S. 425 (1977)\par}
\vspace{0.3em}
{\normalsize Case Brief --- One-Page Edition\par}
\end{center}

\textbf{Memory Jogger.} Seizing and screening Nixon's papers was no bill of attainder and no separation-of-powers violation (\href{https://supreme.justia.com/cases/federal/us/433/425/}{Justia}).

\textbf{Facts.} After Nixon resigned, a GSA agreement let him destroy his tapes. Congress ordered GSA custody and archival screening of his materials, return of personal items, and compensation for any taking.

\textbf{Procedural History.} A three-judge district court rejected his facial challenge.

\textbf{Votes.} Affirmed 7--2: Brennan for the Court; four concurrences; Burger and Rehnquist dissented.

\textbf{Holding.} Nixon was a ``legitimate class of one,'' the Act was not punishment, and it did not unduly disrupt the Executive.

\textbf{Reasoning.} Only Nixon's papers faced destruction; there was no traditional punishment, and the Act served nonpunitive goals. Custody stayed with executive archivists, two Presidents supported the law, and a former President's privilege carries diminished weight. Stevens stressed the resignation and pardon; Rehnquist warned of a ``veritable sword of Damocles.''

\textbf{Examples.} \emph{Same-Side:} securing one official's server. \emph{Opposite-Side:} uncompensated forfeiture amid condemnation. \emph{Fence-Sitter:} a hostile successor's waiver.

\textbf{Critique.} The ruling made the Attainder Clause deferential and invites pretextual classes of one.

\textbf{Key Quotation.} ``The Court, therefore, often has looked beyond mere historical experience and has applied a functional test of the existence of punishment, analyzing whether the law under challenge, viewed in terms of the type and severity of burdens imposed, reasonably can be said to further nonpunitive legislative purposes.'' --- Justice Brennan, 433 U.S. at 475--476.

\end{document}
